\hypertarget{appendix-3999c6209649}{}
\section*{Saturating arithmetic}
\section{Saturating arithmetic}\label{saturating-arithmetic}

\textbf{Curriculum depth:} Introduced · \textbf{First seen:} C002 · \textbf{Latest development:} C002

\subsection{Definition and mechanism}\label{definition-and-mechanism}

Clamp unrepresentable results to range boundaries instead of wrapping; useful for bounded media values but generally non-associative.

Saturation clamps a mathematical result to a permitted interval: unsigned-byte 250+20 becomes 255 instead of 14. Clamp a widened result. Signed saturation is not associative: saturating (100+100)−100 yields 27 while 100+saturating(100−100) yields 100.

\subsection{Prerequisites}\label{prerequisites}

\begin{itemize}
\tightlist
\item
  \hyperref[appendix-2ed044358b82]{Arithmetic policies} --- A policy states the intended response to an unrepresentable answer: wrap modulo a base, report checked failure, widen storage, clamp to an endpoint, or grow arbitrary precision.
\item
  \hyperref[appendix-1a599c924a4e]{Signed integer overflow} --- Overflow concerns the operation's signed type after promotions.
\end{itemize}

\subsection{Used by}\label{used-by}

Connect later chunks here when they use this concept.

\subsection{Recall and next depth}\label{recall-and-next-depth}

Explain saturating arithmetic without the label, work the example above, then state the assumption whose removal would change the result.

\textbf{Pending:} SIMD saturation, normalized formats, HDR policy, and evaluation-order effects.

\subsection{Source and navigation}\label{source-and-navigation}

\href{Continuous\%20Notes/Computing\%20Foundations\%20-\%20From\%20Bits\%20to\%20C++\%20and\%20Systems\%20APIs\#Arithmetic\%20becomes\%20a\%20policy}{C002 teaching chapter} · \hyperref[appendix-fd090d982a15]{Part definitions and mechanisms} · \href{}{Dependency Map (02-Foundations-II, appendix-4ad47bdd269a)}
